Let $P$ and $Q$ be posets on disjoint ground sets.  Their
\emph{series composition} $P;Q$ keeps the internal relations of $P$ and
$Q$ and adds $p<q$ for every $p\in P$ and $q\in Q$.  Their
\emph{parallel composition} $P\|Q$ keeps the internal relations and adds
no cross-relations.  A finite poset is series-parallel if it is obtained
from one-element posets by iterating these two operations.  Equivalently,
finite series-parallel posets are the $N$-free posets: this is the
Valdes--Tarjan--Lawler characterization applied to the transitive directed
comparability graph of the poset~\cite{ValdesTarjanLawler1982}.

\begin{lemma}[Extension counts]
\label{lem:extension-counts}
If $|P|=p$ and $|Q|=q$, then
\[
        |L(P;Q)|=|L(P)|\,|L(Q)|
\]
and
\[
        |L(P\|Q)|=\binom{p+q}{p}|L(P)|\,|L(Q)|.
\]
\end{lemma}

\begin{proof}
For $P;Q$, every element of $P$ must precede every element of $Q$, so a
linear extension is an extension of $P$ followed by an extension of $Q$.

For $P\|Q$, choose the $p$ positions occupied by the elements of $P$,
fill them by a linear extension of $P$, and fill the remaining positions
by a linear extension of $Q$.
\end{proof}

\begin{lemma}[Series closure]
\label{lem:series-closure}
If $P_1$ and $P_2$ are extension tilers, then $P_1;P_2$ is an extension
tiler.
\end{lemma}

\begin{proof}
Let $|P_1|=p$, $|P_2|=q$, and let $X$ be a label set of size $p+q$.
Given a word $\sigma\in S_X$, let $Y$ be the set of the first $p$ letters
of $\sigma$ and let $Z=X\setminus Y$.

In a linear extension of a series composition, the label set of the first
component is forced to appear before the label set of the second
component.  Therefore $Y$ is forced for any tile containing $\sigma$.
Restrict $\sigma$ to its subwords on $Y$ and on $Z$.  By relabeling
invariance and the extension-tiler property of $P_1$ and $P_2$, these
subwords lie in unique component tiles.  Combining those two component
tiles gives a tile for a relabeling of $P_1;P_2$ containing $\sigma$.

The forced block $Y$ and uniqueness inside the two component tilings give
disjointness.  Hence the constructed family is a partition of $S_X$.
\end{proof}

\begin{lemma}[Parallel closure]
\label{lem:parallel-closure}
If $P_1$ and $P_2$ are extension tilers, then $P_1\|P_2$ is an extension
tiler.
\end{lemma}

\begin{proof}
Let $|P_1|=p$, $|P_2|=q$, and fix one decomposition
$X=Y\sqcup Z$ with $|Y|=p$ and $|Z|=q$.  Given
$\sigma\in S_X$, read the subword $\sigma|_Y$ and the subword
$\sigma|_Z$.  By relabeling the component tilings to the arbitrary label
sets $Y$ and $Z$, these determine unique relabeled tiles for $P_1$ on
$Y$ and for $P_2$ on $Z$.

Since parallel composition adds no cross-relations, every shuffle of a
linear extension from the first component with a linear extension from
the second component is a linear extension of the parallel composition.
Thus the two component tiles determine a tile of $P_1\|P_2$ containing
$\sigma$.  Disjointness follows from the uniqueness of the two subwords'
component tiles.
\end{proof}

\begin{theorem}[Series-parallel extension-tiler theorem]
\label{thm:sp-extension-tiler}
Every finite series-parallel poset is an extension tiler.
\end{theorem}

\begin{proof}
The one-element poset is an extension tiler.  By
Lemmas~\ref{lem:series-closure} and~\ref{lem:parallel-closure}, the class of extension
tilers is closed under series and parallel composition.  Therefore every
poset built from singletons by these operations is an extension tiler.
\end{proof}

\begin{corollary}
If $P$ is series-parallel on $n$ elements, then $|L(P)|$ divides $n!$.
\end{corollary}

\begin{proof}
By Theorem~\ref{thm:sp-extension-tiler}, $S_n$ is a disjoint union of
translates of $L(P)$, so $|L(P)|$ divides $|S_n|=n!$.
\end{proof}

\begin{remark}
The converse divisibility statement is false for $n=5$: there are
non-series-parallel posets on five elements with $|L(P)|$ dividing
$120$.  The next section certifies that these divisible false positives
still do not tile by left translates.
\end{remark}
